%%%%%%%%%%%%%%%%%%%%%%%%%%%%%%%%%%%%%%%%%%%%%%%%%%%%%%%%%%%%%%%%%%%%%%%%%%%%%%%%
% Objetivo: Lista de términos empleados en el documento,                      %
%            junto con sus respectivos significados.                          %
%%%%%%%%%%%%%%%%%%%%%%%%%%%%%%%%%%%%%%%%%%%%%%%%%%%%%%%%%%%%%%%%%%%%%%%%%%%%%%%%

\newglossaryentry{biosenial}{
  name=bioseñal,
  description={Señal eléctrica, magnética o de otro tipo generada por organismos vivos, como las señales EEG del cerebro o EMG de los músculos.}
}

\newglossaryentry{neurofeedback}{
  name=neurofeedback,
  description={Técnica que permite a una persona observar su actividad cerebral en tiempo real para aprender a modificarla voluntariamente.}
}

\newglossaryentry{artefacto}{
  name=artefacto,
  description={Señal no deseada en los datos EEG causada por movimientos oculares, parpadeos, actividad muscular o interferencias eléctricas.}
}

\newglossaryentry{epoch}{
  name=época,
  description={Segmento temporal específico de datos EEG, típicamente sincronizado con un evento o estímulo particular.}
}

\newglossaryentry{montage}{
  name=montaje,
  description={Configuración específica de electrodos EEG que define cómo se combinan o referencian las señales de diferentes canales.}
}

\newglossaryentry{impedancia}{
  name=impedancia,
  description={Resistencia eléctrica entre el electrodo y el cuero cabelludo, que debe mantenerse baja para una buena calidad de señal EEG.}
}

\newglossaryentry{streaming}{
  name=streaming,
  description={Transmisión continua de datos en tiempo real, permitiendo el procesamiento inmediato sin necesidad de almacenamiento previo.}
}

\newglossaryentry{paradigma}{
  name=paradigma BCI,
  description={Método específico para extraer información de señales cerebrales, como imaginación motora, P300 o SSVEP.}
}

\newglossaryentry{topografia}{
  name=topografía,
  description={Representación visual de la actividad eléctrica cerebral distribuida sobre el cuero cabelludo en forma de mapa de colores.}
}

\newglossaryentry{sincronizacion}{
  name=sincronización temporal,
  description={Proceso de alinear temporalmente señales de diferentes fuentes para establecer correlaciones válidas entre eventos cerebrales y acciones del usuario.}
}

\newglossaryentry{dataset}{
  name=dataset,
  description={Conjunto estructurado de datos que contiene señales EEG sincronizadas con eventos de entrada, utilizado para entrenar modelos de aprendizaje automático en aplicaciones BCI.}
}

\newglossaryentry{desincronizacion}{
  name=desincronización relacionada con eventos,
  description={Disminución en la potencia de ciertas bandas de frecuencia EEG que ocurre cuando se planifica o ejecuta un movimiento, especialmente visible en ritmos mu y beta sobre áreas motoras.}
}

\newglossaryentry{rechazo_modo_comun}{
  name=rechazo de modo común,
  description={Capacidad de un amplificador diferencial para rechazar señales que aparecen simultáneamente en ambas entradas (como interferencias eléctricas), mientras amplifica las diferencias de voltaje entre ellas.}
}

\newglossaryentry{boilerplate}{
  name=código boilerplate,
  description={Código repetitivo y estándar que debe escribirse en muchos lugares con poca alteración, como plantillas o estructuras básicas que se repiten frecuentemente en la programación.}
}

\newglossaryentry{insights}{
  name=insights,
  description={Comprensiones profundas o percepciones valiosas obtenidas del análisis de datos, que revelan patrones, tendencias o relaciones no evidentes a primera vista.}
}

\newglossaryentry{footprint}{
  name=footprint,
  description={Huella o espacio ocupado por un software en términos de memoria, almacenamiento o recursos del sistema necesarios para su funcionamiento.}
}

\newglossaryentry{callback}{
  name=callback,
  description={Función que se ejecuta como respuesta a un evento específico, permitiendo procesamiento asíncrono sin bloquear la ejecución principal del programa.}
}

\newglossaryentry{polling}{
  name=polling,
  description={Técnica de programación donde un programa verifica repetidamente el estado de un recurso o dispositivo en intervalos regulares para detectar cambios o nuevos datos.}
}